\documentclass[10pt,a4paper]{amsart}
\usepackage[margin=19mm]{geometry}
\usepackage{amsmath,amssymb,mathtools}
\usepackage[scr=rsfs,cal=euler]{mathalfa}
\usepackage{xfrac}
\usepackage[unicode,hidelinks,bookmarks=false]{hyperref}
\newcommand{\ie}{i.e.\ }
\newcommand{\cf}{cf.\ }
\newcommand{\resp}{respectively\ }
\newcommand{\Subsection}{\subsection}
\providecommand{\nobreakdash}{\nobreak\hbox{-}}
\setlength{\emergencystretch}{2em}
\multlinegap0pt
\usepackage{bm}
\usepackage{bbm}
\usepackage{dsfont}
\usepackage{xspace}
\usepackage{cancel}
\usepackage{mathtools}
\usepackage{amsthm}
\makeatletter
\@ifundefined{lemma}{\newtheorem{lemma}{Lemma}}{}
\makeatother
\title[Global Attractors of Non-autonomous Lattice Dynamical System]{Global Attractors of Non-autonomous Lattice Dynamical Systems}
\author{David Cheban, Andrei Sultan}
\date{Source: arXiv:2506.17841v1 (CC BY 4.0)}
\begin{document}
\maketitle

\noindent\emph{Source and licence.} Global Attractors of Non-autonomous Lattice Dynamical Systems. Version arXiv:2506.17841v1 (CC BY 4.0), distributed under the Creative Commons Attribution 4.0 International licence, as recorded in the arXiv licence metadata for this exact version and shown on the abstract page https://arxiv.org/abs/2506.17841v1. Full rights notice: licenses/arxiv-2506.17841-CC-BY-4.0.txt.\par\smallskip

\noindent\textit{Editorial scope.} Counted unit: the Lemma whose source label is \texttt{lCD1} in the article (source0.tex lines 634--651, source label \texttt{lCD1}), delivered together with the article's own complete original proof of it (source0.tex lines 652--767, one \texttt{proof} environment of 116 lines). No mathematics is written, repaired or re-derived: statement and proof are the article's own text line for line. Required context retained uncounted: none. Every definition, display and label of those lines is retained unchanged. Omitted from this package and not redistributed: 1--633, 768--1247. Nothing inside a retained excerpt is omitted or altered: every retained body is delivered exactly as the article's lines are written, so no omission receipt of any kind accompanies this package. Citations: the retained excerpts carry no citation command at all, so no bibliography entry of the article is transcribed and no reference is redistributed. Reference adaptation: none was needed and none was made; every reference inside the delivered unit resolves inside it, so no unresolved reference or citation is produced. Printed slips kept literal: no printed word, symbol, hypothesis or number of the counted statement or of its proof was altered. Layout and class: the article's own class is \texttt{amsart} with packages including amsmath, amsfonts, amssymb, graphicx, xcolor, hyperref; the wrapper is the lane's pinned \texttt{amsart} editorial wrapper at 10pt on A4, which restates the theorem-like environments the retained text needs and adds the article's own macro definitions that the retained text needs (none), with extra wrapper packages (bm, bbm, dsfont, xspace, cancel, mathtools). The delivered numbering is the unit's own. Assets: no retained span contains a figure environment, a graphics-inclusion command, a PDF-page inclusion, a table or a TikZ picture; no image file is copied into this package, the package declares no asset dependency and no image file is redistributed with it. Licence: version arXiv:2506.17841v1 of this article is distributed under the Creative Commons Attribution 4.0 International licence, as recorded in the arXiv licence metadata for this exact version and shown on the abstract page https://arxiv.org/abs/2506.17841v1. A full rights notice is delivered at licenses/arxiv-2506.17841-CC-BY-4.0.txt. Characters: every retained excerpt is pure ASCII, so the delivered engine, XeLaTeX with its default Latin Modern text font, reports no missing character.
\begin{lemma}\label{lCD1} Assume that the skew product dynamical
system $\langle \ell_{2},\varphi,(Y,\mathbb T,\sigma)\rangle$
satisfies the following conditions:
\begin{enumerate}
\item the metric space $Y$ is compact; \item the cocycle $\varphi$
admits a bounded absorbing set $K\subset \ell_{2}$, i.e., for any
bounded subset $B\subset \ell_{2}$ there exists a positive number
$L=L(B)$ such that $\varphi(t,B,Y)$ $\subset $ $K$ for any $t\ge
L$ (or equivalently: $\mathcal K:=K\times Y$ is an absorbing set
for the skew-product dynamical system $(X,\mathbb R_{+},\pi)$
generated by cocycle $\varphi$); \item the cocycle $\varphi$
satisfy an asymptotic tails property on the absorbing set
$\mathcal K \subset X$.
\end{enumerate}

Then the skew-product dynamical system $(X,\mathbb R_{+},\pi)$
generated by the cocycle $\varphi$ is asymptotically compact.
\end{lemma}

\begin{proof} Let $B$ be a bounded, closed and positively invariant
subset of $X$. Since the space $Y$ is compact then there exists a
positive number $r>0$ such that $B\subseteq B[0,r]\times Y$.
Consider the sequences $\{x_n\}\subset B$ and $t_n\to +\infty$ as
$n\to \infty$. We will show that the sequence $\{\pi(t_n,x_n)\}$
is precompact in the space $X:=\ell_{2}\times Y$ or equivalently
the sequence $\{\varphi(t_n,u_n,y_n)\}$ is precompact in
$\ell_{2}$ ($x_n=(u_n,y_n)$ and
$\pi(t_n,x_n)=(\varphi(t_n,u_n,y_n),\sigma(t_n,y_n))$ for any
$n\in \mathbb N$) because $\{y_n\}\subset Y$ and the space $Y$ is
compact.

Since the est $B$ is positively invariant then $\pi(t,B)\subseteq
B\subseteq B[0,r]\times Y$ for any $t\ge 0$ and, consequently,
\begin{equation}\label{eqBK1}
\varphi(t,u,y)\in B[0,r] \nonumber
\end{equation}
for any $(u,y)\in B$ and $t\ge 0$. In particular, we have
\begin{equation}\label{BK2}
|\varphi(t_n,u_n,y_n)|\le r \nonumber
\end{equation}
for any $n\in \mathbb N$). Thus the sequence
$\{\varphi(t_n,u_n,y_n)\}$ is weakly precompact in $\ell_{2}$ and,
consequently, we can assume that $\{\varphi(t_n,u_n,y_n\}$ is
weakly convergent, i.e,
\begin{equation}\label{eqWC1}
\varphi(t_n,u_n,y_n)\rightharpoonup u \nonumber
\end{equation}
weakly in $\ell_{2}$.

Let $\varepsilon$ be an arbitrary positive number. Since $K\subset
\ell_{2}$ is an absorbing set for the cocycle $\varphi$ then there
exists a positive number $L_{1}=L_{1}(B[0,r])$ such that
\begin{equation}\label{eqAS_01}
\varphi(t,u,y)\in K
\end{equation}
for any $t\ge L_{1}(B[0,r])$ and $(u,y)\in B[0,r]\times Y$. Let
$n_1\in \mathbb N$ be a number such that $t_n\ge L_{1}$ for any
$n\ge n_1$ and, consequently, from (\ref{eqAS_01}) we obtain
\begin{equation}\label{eqAS_2}
\varphi(t_n,u_n,y_n)\in K
\end{equation}
for any $n\ge n_1$.

We will show that the sequence $\{\varphi(t_n,u_n,y_n)\}$ converge
in the space $\ell_{2}$, that is, for any $\varepsilon >0$ there
exists a number $n(\varepsilon)\in \mathbb N$ such that
\begin{equation}\label{eqAS3}
\|\varphi(t_n,u_n,y_n)-u\|<\varepsilon \nonumber
\end{equation}
for any $n\ge n(\varepsilon)$.

By (\ref{eqAS_2}) we have
\begin{equation}\label{eqAS4}
\varphi(L_{1},u_n,y_n)\in K
\end{equation}
for any $n\in \mathbb N$. Since the cocycle $\varphi$ satisfy an
asymptotic tails on the set $K$ for given $\varepsilon
>0$ there exist $k_{1}(\varepsilon)\in \mathbb N$ and
$L_{1}(\varepsilon)>0$ such that
\begin{equation}\label{eqAS5}
\sum\limits_{|i|\ge
k_{1}(\varepsilon)}|\varphi_{i}(t,\varphi(L_{1},u_n,y_n),\sigma(L_{1},y_n))|^{2}<\varepsilon^{2}/8
\end{equation}
for any $t\ge L_{1}(\varepsilon)$.

Since $t_n\to +\infty$ as $n\to \infty$, there exists
$n_{2}(\varepsilon)\in \mathbb N$ such  that if $n\ge
n_{2}(\varepsilon)$, then $t_n - L_{1}\ge L_{1}(\varepsilon)$, and
hence from (\ref{eqAS5}), we have
\begin{eqnarray}\label{eqAS6}
& \sum\limits_{|i|\ge
k_{1}(\varepsilon)}|\varphi_{i}(t_n,u_n,y_n)|^{2}=
\\
& \sum\limits_{|i|\ge
k_{1}(\varepsilon)}|\varphi_{i}(t_n-L_{1},\varphi(L_{1}u_n,y_n),\sigma(L_1,y_n))|^{2}\le
\varepsilon^{2}/8.\nonumber
\end{eqnarray}

Since $u\in \ell_{2}$ then there exists $k_{2}(\varepsilon)$ such
that
\begin{equation}\label{eqAS7}
\sum\limits_{|i|\ge k_{2}(\varepsilon)}|u_{i}|^{2}\le
\varepsilon^{2}/8.
\end{equation}
Let $k(\varepsilon)=\max\{k_1(\varepsilon),k_{2}(\varepsilon)\}$.
By the weak convergence of $\{\varphi(t_n,u_n,y_n)\}$ we have that
\begin{equation}\label{eqAS8_1}
\varphi_{i}(t_n,u_n,y_n)\to u_{i}\nonumber
\end{equation}
as $n\to \infty$ (for any $|i|\le k(\varepsilon)$), which implies
that there exists $n_{3}(\varepsilon)$ such that for any $n\ge
n_{3}(\varepsilon)$ we obtain
\begin{equation}\label{eqAS8}
\sum\limits_{|i|\le
k(\varepsilon)}|\varphi_{i}(t_n,u_n,y_n)-u_{i}|^{2}\le
\varepsilon^{2}/2
\end{equation}
for any $n\ge n_{3}(\varepsilon).$

Setting
$n(\varepsilon):=\max\{n_{1},n_{2}(\varepsilon),n_{3}(\varepsilon)\}$,
from (\ref{eqAS6})-(\ref{eqAS8}) we get that, for $n\ge
n(\varepsilon)$,
\begin{eqnarray}\label{eqAS9}
& \|\varphi(t_n,u_n,y_n)-u\|^{2}=\sum\limits_{|i|\le
k(\varepsilon)}|\varphi_{i}(t_n,u_n,y_n)-u_{i}|^{2} +\nonumber\\
& \sum\limits_{|i|>
k(\varepsilon)}|\varphi_{i}(t_n,u_n,y_n)-u_{i}|^{2}\le
\varepsilon^{2}/2+2\sum\limits_{|i|\ge
k_{\varepsilon}}(|\varphi_{i}(t_n,u_n,y_n)|^{2}+|u_{i}|^{2})\le
\varepsilon^{2}.\nonumber
\end{eqnarray}
as desired. Hence we obtain $\varphi(t_n,u_n,y_n)$ converges to $u$
in the space $\ell_{2}$. Lemma is completely proved.
\end{proof}

\end{document}
